\begin{trapbox}
\begin{center}
\begin{minipage}{0.98\linewidth}
\centering
\textbf{同一曲线，限制定义域会改变最小值}\par\medskip
\begin{tikzpicture}[font=\small,>=Stealth]

\begin{scope}[x=2cm,y=1.05cm]
\draw[->] (0,0)--(3.7,0) node[right] {$x$};
\draw[->] (0,0)--(0,4.5) node[above] {$y$};
\draw[gray,dashed,thick,domain=0.26:3.4,samples=90] plot (\x,{\x+1/\x});
\draw[CTrap,very thick,domain=2:3.4,samples=60] plot (\x,{\x+1/\x});
\draw[dotted] (1,0)--(1,2) (2,0)--(2,2.5)--(0,2.5);
\draw[gray,fill=white] (1,2) circle (2.5pt);
\fill[CTrap] (2,2.5) circle (2.5pt);
\node[below] at (1,0) {$1$};\node[below] at (2,0) {$2$};
\node[fill=white] at (1,1.25) {原最低点被排除};
\node[above right,fill=white] at (2,2.5) {新最低点 $(2,5/2)$};
\node[left] at (0,2.5) {$5/2$};
\end{scope}
\end{tikzpicture}\par\smallskip
灰虚线：$x>0$，值域 $[2,+\infty)$；强调实线：$x\ge2$，值域 $[5/2,+\infty)$。这里 $x=1$ 不可取，新端点 $x=2$ 可取。
\end{minipage}
\end{center}


	\begin{description}[style=nextline, leftmargin=2.8em, labelsep=0.5em, itemsep=0.35em]
		\item[\textbf{\textcolor{CTrap}{易错点一（正负区间混淆）}}] 误以为 $x<0$ 时也有 $f(x)\ge2\sqrt{AB}$，将最大值当成最小值.
		\item[\textbf{\textcolor{CTrap}{正确理解（不等号方向）}}] $x<0$ 时，$Ax<0,\frac{B}{x}<0$，利用基本不等式时需提取负号，得到 $f(x)\le-2\sqrt{AB}$，此为最大值.
		\item[\textbf{\textcolor{CTrap}{错因分析（忽视符号）}}] 受正数区间思维定势，未考虑正负改变不等号方向.
	\end{description}

	\medskip
	\begin{description}[style=nextline, leftmargin=2.8em, labelsep=0.5em, itemsep=0.35em]
		\item[\textbf{\textcolor{CTrap}{易错点二（换元定义域忽略）}}] 换元后直接使用对勾结论，不考虑新变量的正负.
		\item[\textbf{\textcolor{CTrap}{正确理解（先定范围）}}] 对 $t$ 的范围需分类：$t>0$ 时用 $f(t)\ge2\sqrt{AB}$，$t<0$ 时用 $f(t)\le-2\sqrt{AB}$.
		\item[\textbf{\textcolor{CTrap}{错因分析（缺乏分类意识）}}] 换元后未养成先求定义域再套用公式的习惯.
	\end{description}

	\medskip
	\begin{description}[style=nextline, leftmargin=2.8em, labelsep=0.5em, itemsep=0.35em]
		\item[\textbf{\textcolor{CTrap}{易错点三（等号条件缺失）}}] 求出最值 $2\sqrt{AB}$ 后直接写进值域，不检验是否取得到.
		\item[\textbf{\textcolor{CTrap}{正确理解（验证取等）}}] 必须检查 $x=\pm\sqrt{\frac{B}{A}}$ 是否在定义域内；若取等点被限制条件排除，原来的最值不再能直接照搬，须在新定义域内重新求值域，并判断端点能否取到.
		\item[\textbf{\textcolor{CTrap}{错因分析（忽视取等条件）}}] 只记住结论形式，未理解等号成立需要 $Ax=\frac{B}{x}$ 有解.
	\end{description}


\medskip
\begin{center}
\begin{minipage}{0.98\linewidth}
\centering
\textbf{斜线段投到 $y=x$ 方向，长度只会缩短}\par\medskip
\begin{tikzpicture}[font=\small,>=Stealth]

\begin{scope}[x=1.6cm,y=1.6cm]
\draw[->,gray] (-0.3,0)--(3.2,0) node[right] {水平差};
\draw[->,gray] (0,-0.3)--(0,2.8) node[above] {竖直差};
\draw[dashed] (0,0)--(2.6,2.6) node[above right] {$y=x$ 方向};
\draw[CProof,very thick] (0,0)--(3,1) node[right] {$(u,v)$};
\draw[CExample,very thick] (0,0)--(2,2);
\draw[dotted,thick] (3,1)--(2,2);
\draw (2.12,1.88)--(2,1.76)--(1.88,1.88);
\node[below] at (1.6,0.4) {$d$};
\node[above left,rotate=45] at (1.45,1.45) {$(u+v)/\sqrt2$};
\node[below left] at (0,0) {$O$};
\end{scope}
\end{tikzpicture}\par\smallskip
平移线段使起点到原点，图示 $u=3,v=1$。投影终点为 $((u+v)/2,(u+v)/2)$；只有 $u=v$ 时，投影长等于 $d$。
\end{minipage}
\end{center}

\begin{description}[style=nextline, leftmargin=2.8em, labelsep=0.5em, itemsep=0.4em]
\item[\textbf{\textcolor{CTrap}{易错点四（反比例函数与对勾函数混为一谈）}}]
误认为 $y=\frac{k}{x}$ 本身就是 $Ax+\frac Bx$（$A,B>0$）。正确说法是：在曲线上取点时，坐标和 $x+y=x+\frac{k}{x}$ 才是可用的对勾函数。
\item[\textbf{\textcolor{CTrap}{易错点五（把坐标和当成线段长度）}}]
$\Delta x+\Delta y$ 不是两点间距离，$\frac{\Delta x+\Delta y}{\sqrt2}$ 也只是距离的下界；只有 $\Delta x=\Delta y$ 时，后者才等于实际距离。必须把勾股不等式和取等条件写完整。
\item[\textbf{\textcolor{CTrap}{易错点六（对称图形不等于任意点都关于原点对称）}}]
题意只保证 $P$、$Q$ 分别位于两支曲线上，不能一开始就未经证明地设 $Q=-P$。可先设 $P(a,k/a)$、$Q(-b,-k/b)$（$a,b>0$），证明最小值后再发现最优两点关于原点对称。
\item[\textbf{\textcolor{CTrap}{易错点七（忽略参数及额外限制）}}]
结论 $2\sqrt{2k}$ 只适用于 $k>0$ 且两点能在两条完整分支上自由选取。若限制点的横坐标或可选范围，原取等点可能被排除；若 $k<0$，曲线位于第二、第四象限，题设“第一、第三象限各一点”就不成立。
\end{description}

\end{trapbox}
